\subsection{紧开拓扑}

定理 2. 设 X 是拓扑空间，Y 是一致空间。对 X 的每个紧子集 K 及 Y 的每个开子集 U，用 \(\mathbf{T}(\mathbf{K},\mathbf{U})\) 表示一切连续映射 \(u:X\to Y\) （满足 \(u(\mathbf{K})\subset U\) ）所成之集。则诸集 \(\mathbf{T}(\mathbf{K},\mathbf{U})\) 生成 \(\mathcal{C}(\mathbf{X};\mathbf{Y})\) 上的紧收敛拓扑。

设 \(Y'\) 是与 \(Y\) 相伴的豪斯多夫一致空间（第 II 章，§ 3，第 8 号），并设 \(i: Y \to Y'\) 是 \(Y\) 到 \(Y'\) 上的典范映射。紧收敛拓扑是使映射 \(u \to (i \circ u)|K\) （由 \(\mathcal{C}(X; Y)\) 到 \(\mathcal{C}_u(K; Y')\) 中）连续的最粗拓扑，其中 \(K\) 取遍 \(X\) 的一切紧子集所成之集（第 1 节，第 4 号，命题 4）。因此，对 \(\mathcal{C}_c(X; Y)\) 的拓扑，取 \(\mathcal{C}_c(K; Y')\) 的拓扑的一个子基（对每个紧子集 \(K\) ，在 \(X\) 中），再取这些子基的逆像之并{[}在 \(\mathfrak{P}(\mathcal{C}(X; Y))\) 中{]}，逆像取在 \(\mathcal{C}(X, Y)\) 中，便得到它的一个子基。另一方面，\(Y\) 的每个开子集形如 \(\overline{i}^1(U')\) ，其中 \(U'\) 在 \(Y'\) 中开（第 II 章，§ 3，第 7 号，命题 12）；因此，对每个紧子集 \(K' \supset K\) ，\(T(K, \overline{i}^1(U'))\) 是 \(T(K, U')\) 在映射

\[
\mathcal {C} (\mathbf {X}; \mathbf {Y}) \rightarrow \mathcal {C} _ {u} (\mathbf {K} ^ {\prime}; \mathbf {Y} ^ {\prime}).
\]

于是只剩下来证明定理在 X 紧且 Y 豪斯多夫的情形；以下我们就作这些假设。

首先证明 \(T(\mathbf{K}, \mathbf{U})\) 在 \(\mathcal{C}_c(\mathbf{X}; \mathbf{Y})\) 中是开的。设 \(u_0\) 是这个集的一个点；由于 \(u_0(\mathbf{K})\) 紧（第 I 章，§ 9，第 4 号，定理 2 的推论 1）且含于开集 \(\mathbf{U}\) ，存在对称围域 \(\mathbf{V}\) （\(\mathbf{Y}\) 的围域）使 \(\mathbf{V}(u_0(\mathbf{K})) \subset \mathbf{U}\) （第 II 章，§ 4，第 3 号，命题 4 的推论）。设 \(\mathbf{W}\) 是 \(u_0\) 在 \(\mathcal{C}_c(\mathbf{X}; \mathbf{Y})\) 中的由一切连续映射 \(u: \mathbf{X} \to \mathbf{Y}\) （满足 \((u(x), u_0(x)) \in \mathbf{V}\) 对一切 \(x \in \mathbf{K}\) ）组成的邻域。对这些映射我们显然有 \(u(\mathbf{K}) \subset \mathbf{V}(u_0(\mathbf{K})) \subset \mathbf{U}\) ；因此 \(u \in T(\mathbf{K}, \mathbf{U})\) ，从而 \(W \subset T(\mathbf{K}, \mathbf{U})\) ，这就证明了我们的断言。

反之，若 W 是一点 \(u_{0} \in \mathcal{C}_{c}(\mathbf{X}; \mathbf{Y})\) 的邻域，我们要证明 W 含有有限个形如 \(T(\mathbf{K}, \mathbf{U})\) 的邻域的交。可以设 W 是一切 \(u \in \mathcal{C}(\mathbf{X}; \mathbf{Y})\) （满足 \((u(x), u_{0}(x)) \in \mathbf{V}\) 对一切 \(x \in X\) ）所成之集，其中 V 是 Y 的一个给定围域。由于 \(u_{0}\) 在 X 上连续，它一致连续（第 II 章，§ 4，第 1 号，定理 2）。设 \(V_{1}\) 是 Y 的一个对称围域，在 \(Y \times Y\) 中开，且 \(V_{1}^{2} \subset V\) 。X 可被有限个紧集 \(K_{i}\) （ \(1 \leqslant i \leqslant n\) ）覆盖，使得每个 \(u_{0}(\mathbf{K}_{i})\) 是 \(V_{1}\) -小的（ \(1 \leqslant i \leqslant n\) ）。设 \(U_{i}\) 是开集 \(V_{1}(u_{0}(\mathbf{K}_{i}))\) ，并设 \(u: X \to Y\) 是含于 n 个集 \(T(\mathbf{K}_{i}, \mathbf{U}_{i})\) （它们是 \(u_{0}\) 的邻域）之交的一个连续映射。那么，对每个 \(x \in K_{i}\) ，\(u(x)\) 属于 \(U_{i}\) ，因此 \(u_{0}(x)\) 与 \(u(x)\) 是 \(V_{1}^{2}\) -接近的，从而是 V-接近的。由于每个 \(x \in X\) 属于某个 \(K_{i}\) ，我们有 \(u \in W\) ，证明完成。

这个结果引导我们作出如下的定义：

定义 1. 设 X、Y 是两个拓扑空间，不必可一致化。对 X 的每个紧子集 K 及 Y 的每个开子集 U，设 T(K, U) 是一切 \(u \in \mathcal{C}(X; Y)\) （满足 \(u(K) \subset U\) ）所成之集。C(X; Y) 上由诸集 T(K, U) 生成的拓扑称为紧收敛拓扑或紧开拓扑；我们把赋有这个拓扑的 C(X; Y) 所得到的拓扑空间记作 \(\mathcal{C}_{c}(X; Y)\) 。

若 Y 是一致空间，则由定理 2，这个定义与第 1 节第 3 号中给出的定义一致。

若 H 是 \(\mathcal{C}(\mathbf{X};\mathbf{Y})\) 的一个子集，我们把由 \(\mathcal{C}_{c}(\mathbf{X};\mathbf{Y})\) 的拓扑在 H 上诱导的拓扑称为 H 上的紧开拓扑。

例。设 I 是 R 中的区间 \([0, 1]\) 。若 Y 是任一拓扑空间，空间 \(\mathcal{C}_{c}(I; Y)\) 称为 Y 中的道路空间。对每个 \(y \in Y\) ，子空间 \(\Omega_{y}(Y)\) ——\(\mathcal{C}_{c}(I; Y)\) 中由满足 \(u(0) = u(1) = y\) 的道路 u 组成的子空间——称为（Y 中）点 y 处的闭路空间。

注。1) 同样，\(\mathcal{C}(\mathbf{X};\mathbf{Y})\) 上由 \(\mathbf{Y}^{\mathbf{X}} = \mathcal{F}(\mathbf{X};\mathbf{Y})\) 上的乘积拓扑诱导的拓扑称为逐点收敛拓扑（Y 不必可一致化）；它由形如 \(T(\{x\}, \mathbf{U})\) 的集生成，其中 \(x\) 取遍 \(\mathbf{X}\) ，\(\mathbf{U}\) 取遍 \(\mathbf{Y}\) 的一切开子集所成之集，因此它比紧开拓扑粗。由此推出，若 \(\mathbf{Y}\) 豪斯多夫，则空间 \(\mathcal{C}_s(\mathbf{X}; \mathbf{Y})\) 豪斯多夫（第 I 章，§ 8，第 1 号，命题 5 的推论）。

\begin{enumerate}
\def\labelenumi{\arabic{enumi})}
\setcounter{enumi}{1}
\tightlist
\item
  设 \(\mathfrak{S}\) 是 Y 的拓扑的一个子基，并设 \(\mathfrak{R}\) 是 X 的一些紧子集所成之集，具有下列性质：
\end{enumerate}

\begin{enumerate}
\def\labelenumi{(\Alph{enumi})}
\setcounter{enumi}{17}
\tightlist
\item
  若 \(\mathbf{L}\) 是 \(\mathbf{X}\) 的任一紧子集且 \(\mathbf{V}\) 是 \(\mathbf{L}\) 的任一邻域，则存在有限个集 \(\mathbf{K}_i \in \mathfrak{R}\) 使 \(\mathbf{L} \subset \bigcup_{i} \mathbf{K}_i \subset \mathbf{V}\) 。
\end{enumerate}

那么诸集 \(T(\mathbf{K}, \mathbf{U})\) （其中 \(\mathbf{K} \in \mathfrak{R}\) 且 \(\mathbf{U} \in \mathfrak{S}\) ）构成 \(\mathcal{C}(\mathbf{X}; \mathbf{Y})\) 上紧开拓扑的一个子基。为证明这一点，我们必须证明：若 \(\mathbf{L}\) 是 \(\mathbf{X}\) 的任一紧子集且 \(\mathbf{V}\) 是 \(\mathbf{Y}\) 的任一开子集，且 \(u \in T(\mathbf{L}, \mathbf{V})\) ，则存在有限对 \((\mathbf{K}_i, \mathbf{U}_i)\) 使 \(\mathbf{K}_i \in \mathfrak{R}\) 、\(\mathbf{U}_i \in \mathfrak{S}\) 且 \(u \in \bigcap_{i} T(\mathbf{K}_i, \mathbf{U}_i) \subset T(\mathbf{L}, \mathbf{V})\) 。首先注意，对每个有限序列 \((s_k)\) （由 \(\mathfrak{S}\) 中的集合组成）及每个紧子集 \(\mathbf{M}\) （\(\mathbf{X}\) 的子集），由定义我们有 \(T\left(\mathbf{M}, \bigcap_k \mathbf{S}_k\right) = \bigcap_k T(\mathbf{M}, \mathbf{S}_k)\) 。因此我们首先可把 \(\mathfrak{S}\) 换成 \(\mathfrak{S}\) 中集合的有限交所成之集，即我们可设 \(\mathfrak{S}\) 是 \(\mathbf{Y}\) 的拓扑的一个基。由假设，\(u(\mathbf{L})\) 拟紧且含于 \(\mathbf{V}\) ，故存在有限个集 \(\mathbf{U}_i \in \mathfrak{S}\) （含于 \(\mathbf{V}\) ）覆盖 \(u(\mathbf{L})\) 。诸集 \(\overline{u}^1(\mathbf{U}_i)\) 在 \(\mathbf{X}\) 中开且覆盖 \(\mathbf{L}\) 。因此对每个 \(x \in \mathbf{L}\) ，存在紧邻域 \(N_x\) （\(x\) 在 \(\mathbf{L}\) 中的邻域），含于某个 \(\overline{u}^1(\mathbf{U}_i)\) 。\(\mathbf{L}\) 可被有限个这样的集 \(N_{x_j} = L_j\) 覆盖；对每个 \(j\) ，用 \(i(j)\) 表示指标 \(i\) （满足 \(L_j \subset \overline{u}^1(\mathbf{U}_i)\) ）之一。这样，对每个指标 \(j\) 存在{[}由 (R){]}有限个集 \(K_{jk} \subset \overline{u}^1(\mathbf{U}_{i(j)})\) （属于 \(\mathfrak{R}\) ）覆盖 \(L_j\) 。对每个 \(v \in \bigcap_{j,k} T(\mathrm{K}_{jk}, \mathrm{U}_{i(j)})\) 我们有 \(\bigcup_{k} v(\mathrm{K}_{jk}) \subset U_{i(j)}\) ，因此 \(v(\mathrm{L}_j) \subset U_{i(j)}\) ，且 \(v(\mathrm{L}) = \bigcup_{j} v(\mathrm{L}_j) \subset \bigcup_{j} U_{i(j)} \subset V\) ；于是我们的断言得证。

定理 3. 设 X、Y、Z 是三个拓扑空间，并设 f 是 X × Y 到 Z 中的一个映射。若 f 连续，则 \(f: x \to f(x, .)\) 是 X 到 \(\mathcal{C}_{c}(Y; Z)\) 中的一个连续映射。若 Y 局部紧，则逆命题成立。

设 \(f\) 连续。为证明 \(\tilde{f}\) 连续，我们必须证明：对每个紧子集 \(K\) （在 \(Y\) 中）及每个开子集 \(U\) （在 \(Z\) 中），\(V\) （\(T(K, U)\) 在 \(\tilde{f}\) 下的逆像）在 \(X\) 中是开的。设 \(x_0 \in V\) ；对每个 \(y \in K\) ，我们有 \(f(x_0, y) \in U\) ，且由于 \(f\) 连续，存在邻域 \(V_y\) （\(x_0\) 在 \(X\) 中的邻域）及邻域 \(W_y\) （\(y\) 在 \(Y\) 中的邻域）使 \(f(V_y \times W_y) \subset U\) 。由于 \(K\) 紧，存在有限个点 \(y_{i} \in K\) 使诸集 \(W_{y_{i}}(1 \leqslant i \leqslant n)\) 覆盖 K。设 \(V'\) 是诸邻域 \(V_{y_{i}}\) （\(x_{0}\) 的邻域）的交，它是 \(x_{0}\) 的一个邻域；若 \(x \in V'\) 且 \(y \in K\) ，我们有 \(f(x, y) \in U\) ，因为 y 含于某个 \(W_{y_{i}}\) 且 x 含于每个 \(V_{y_{i}}\) ；因此 \(V' \subset V\) ，从而 V 是它的每一点的邻域，于是在 X 中开。

反之，设 \(\tilde{f}\) 连续且 Y 局部紧，我们来证明 \(f\) 连续。设 \(x_0 \in \mathbf{X}\) ，设 \(y_0 \in \mathbf{Y}\) ，并设 U 是 \(f(x_0, y_0)\) 在 Z 中的一个开邻域；我们将证明存在邻域 V （\(x_0\) 在 X 中的邻域）及邻域 W （\(y_0\) 在 Y 中的邻域）使 \(f(\mathbf{V} \times \mathbf{W}) \subset \mathbf{U}\) 。由于 \(y \to f(x_0, y)\) 连续，存在紧邻域 W （\(y_0\) 的邻域）使 \(f(\{x_0\} \times \mathbf{W}) \subset \mathbf{U}\) 。另一方面，由于 \(\tilde{f}\) 连续，集 V （由 \(x \in \mathbf{X}\) 组成，满足 \(f(x, .) \in T(\mathbf{W}, \mathbf{U})\) ）{[}即使 \(f(x, y) \in \mathbf{U}\) 对一切 \(y \in \mathbf{W}]\)是 X 的一个开子集，因而是 \(x_0\) 的一个邻域；且我们有 \(f(\mathbf{V} \times \mathbf{W}) \subset \mathbf{U}\) 。

证毕。

推论 1. 设 X 是局部紧空间，Y 是拓扑空间，H 是 C (X; Y) 的一个子集。则 H 上的紧开拓扑是使映射 \((u, x) \to u(x)\) （由 \(H \times X\) 到 Y 中）连续的最粗拓扑。

因为由定理 3，这个映射连续当且仅当典范单射 \(\mathbf{H} \to \mathcal{C}_c(\mathbf{X}; \mathbf{Y})\) 连续。

注。3) 设 X 是局部紧空间，Y 是豪斯多夫拓扑空间。若 \(\mathcal{G}\) 是 \(\mathcal{C}(\mathbf{X};\mathbf{Y})\) 的一个子集 H 上的一个拓扑，使得映射 \((u,x)\to u(x)\) 在 \(\mathrm{H}\times \mathrm{X}\) 上连续，且 H 关于 \(\mathcal{G}\) 紧，则 \(\mathcal{G}\) 是紧开拓扑。因为由推论 1 它比后者细，且由于紧开拓扑是豪斯多夫的，这两个拓扑相同。注意，若此外 Y 完全正则，则 H 关于与 Y 的拓扑相容的每个一致结构都等度连续（第 2 节，第 5 号，定理 2 的推论 3），且对 X 的每个紧子集 K，集

\[
\mathrm{H} (\mathrm{K}) = \bigcup_ {x \in \mathrm{K}} \mathrm{H} (x)
\]

是紧的，因为它是 \(\mathbf{H} \times \mathbf{K}\) 在连续映射 \((u, x) \to u(x)\) 下的像。

推论 2. 设 X、Y、Z 是三个拓扑空间，使得 X 豪斯多夫且 Y 局部紧。则限制到 \(\mathcal{C}(\mathbf{X} \times \mathbf{Y}; \mathbf{Z})\) 上的典范双射 \(\mathcal{F}(\mathbf{X} \times \mathbf{Y}; \mathbf{Z}) \to \mathcal{F}(\mathbf{X}; \mathcal{F}(\mathbf{Y}; \mathbf{Z}))\) （集合论，R，§4，第 14 号）是 \(\mathcal{C}_c(\mathbf{X} \times \mathbf{Y}; \mathbf{Z})\) 到 \(\mathcal{C}_c(\mathbf{X}; \mathcal{C}_c(\mathbf{Y}; \mathbf{Z}))\) 上的一个同胚。

这个限制当然是一个双射

\[
\rho : \mathcal {C} (\mathrm{X} \times \mathrm{Y}; \mathrm{Z}) \rightarrow \mathcal {C} (\mathrm{X}; \mathcal {C} _ {c} (\mathrm{Y}; \mathrm{Z}))
\]

由定理 3；因此只剩下证明 \(\mathcal{C}(\mathbf{X}\times\mathbf{Y};\mathbf{Z})\) 上的紧开拓扑是 \(\rho\) 下 \(\mathcal{C}(\mathbf{X};\mathcal{C}_{c}(\mathbf{Y};\mathbf{Z}))\) 上紧开拓扑的逆像。由于诸集 \(T(\mathbf{K},\mathbf{U})\) （其中 \(\mathbf{K}\) 是 \(\mathbf{Y}\) 的紧子集且 \(\mathbf{U}\) 是 \(\mathbf{Z}\) 的开子集）构成 \(\mathcal{C}_{c}(\mathbf{Y};\mathbf{Z})\) 的拓扑的一个子基，由注 2 推出，\(\mathcal{C}_{c}(\mathbf{X};\mathcal{C}_{c}(\mathbf{Y};\mathbf{Z}))\) 的拓扑由形如 \(T(J,T(\mathbf{K},\mathbf{U}))\) 的集生成，其中 \(\mathbf{K}\) 与 \(\mathbf{U}\) 如上且 \(J\) 是 \(\mathbf{X}\) 的紧子集。而 \(T(J,T(\mathbf{K},\mathbf{U}))\) 在 \(\rho\) 下的像正是 \(T(J\times K,\mathbf{U})\) ，因此是一个开集；于是我们证明了 \(\rho\) 连续。为证明 \(\rho\) 是一个同胚，我们首先注意形如 \(J\times K\) 的集（在 \(\mathbf{X}\times\mathbf{Y}\) 中）（其中 \(J\) 是 \(\mathbf{X}\) 的紧子集，\(\mathbf{K}\) 是 \(\mathbf{Y}\) 的紧子集）满足注 2 的条件 (R)：因为若 \(\mathbf{L}\) 是 \(\mathbf{X}\times\mathbf{Y}\) 的紧子集且 \(\mathbf{V}\) 是 \(\mathbf{L}\) 在 \(\mathbf{X}\times\mathbf{Y}\) 中的一个邻域，则投影 \(M=pr_{1}(L)\) 、\(N=pr_{2}(L)\) 紧，因为 \(\mathbf{X}\) 与 \(\mathbf{Y}\) 豪斯多夫，且 \(V\cap(M\times N)\) 是 \(\mathbf{L}\) 在紧空间 \(M\times N\) 中的一个邻域，所以 \(\mathbf{L}\) 的每一点在 \(M\times N\) 中有一个形如 \(J\times K\subset V\) （其中 \(J\subset M\) 且 \(K\subset N\) 紧）的邻域；由于 \(\mathbf{L}\) 可被有限个这样的邻域覆盖，断言得证。因此形如 \(T(J\times K;\mathbf{U})\) 的集（其中 \(J\) 是 \(\mathbf{X},\mathbf{K}\) 的紧子集、\(\mathbf{Y}\) 的紧子集，而 \(\mathbf{U}\) 是 \(\mathbf{Z}\) 的开子集）生成 \(\mathcal{C}_{c}(\mathbf{X}\times\mathbf{Y};\mathbf{Z})\) 的拓扑。但我们已经看到 \(T(J\times K,\mathbf{U})\) 在 \(\rho\) 下的像是开集 \(T(J,T(K,\mathbf{U}))\) （在 \(\mathcal{C}_{c}(\mathbf{X};\mathcal{C}_{c}(\mathbf{Y};\mathbf{Z}))\) 中）；因此 \(\rho\) 是一个同胚。

注意，若此外假定 Z 可一致化，则推论 2 是第 1 节第 4 号命题 2 的平凡推论。

命题 9. 设 X、Y、Z 是三个拓扑空间，其中 Y 局部紧。则映射 \((u, v) \to v \circ u\) （由 \(\mathcal{C}_c(X; Y) \times \mathcal{C}_c(Y; Z)\) 到 \(\mathcal{C}_c(X; Z)\) 中）是连续的。

我们必须证明，对 X 的每个紧子集 K 及 Z 的每个开子集 U，由偶 \((u, v)\) （满足 \(v(u(\mathbf{K})) \subset \mathbf{U}\) ）组成的集 R 在 \(\mathcal{C}_{c}(\mathbf{X}; \mathbf{Y}) \times \mathcal{C}_{c}(\mathbf{Y}; \mathbf{Z})\) 中是开的。设 \((u_{0}, v_{0}) \in R\) ；则 \(u_{0}(\mathbf{K})\) 是局部紧空间 Y 的一个紧子集，含于开集 \(\overline{v}_{0}^{1}(\mathbf{U})\) ，因此存在紧邻域 L （\(u_{0}(\mathbf{K})\) 的邻域）含于 \(\overline{v}_{0}^{1}(\mathbf{U})\) （第 I 章，§ 9，第 7 号，命题 10）。集 V （由一切 \(u \in C_{c}(\mathbf{X}; \mathbf{Y})\) （满足 \(u(\mathbf{K}) \subset \mathring{\mathbf{L}}\) ）组成）是 \(u_{0}\) 的一个邻域，而集 W （由一切 \(v \in C_{c}(\mathbf{Y}; \mathbf{Z})\) （满足 \(v(\mathbf{L}) \subset \mathbf{U}\) ）组成）是 \(v_{0}\) 的一个邻域；此外，关系 \((u, v) \in \mathbf{V} \times \mathbf{W}\) 蕴涵 \(v(u(\mathbf{K})) \subset \mathbf{U}\) 。由此即得结论。

